Der Schlitten fährt vor dem Wind mit konstanter Geschwindigkeit. Die Windkraft auf das Segel muss also die Gleitreibungskraft $F_\text{R} = \mu\,m\,g$ ausgleichen, mit $\mu = \muEO$ für Stahl auf Eis (FoTaBe S.~189). Massgebend ist die Geschwindigkeit des Windes relativ zum Segel:
\begin{align}
 v_\text{rel} &= \vREF = \vWE-\vSE = \vRE
\end{align}
Aus
\begin{align}
 \frac{1}{2}\,c_w\,\rho\,A\,v_\text{rel}^2 &= \mu\,m\,g
\end{align}
folgt mit $c_w = \cwPO$ (quadratische Platte, FoTaBe S.~189) und $\rho = \rhLO$ (FoTaBe S.~187)
\begin{align}
 A &= \AEsF \\
   &= \frac{2\cdot\muE\cdot\mEs\cdot\gE}{\cwP\cdot\rhL\cdot\left(\vRE\right)^2} \\
   &= \AEs \approx \result{\AEsP{0}{2}}
\end{align}
Das entspricht einem quadratischen Segel mit der Seitenlänge $\sqrt{A} = \sEs \approx \sEsP{0}{2}$.
